\section{Conclusion}
\label{sec:conclusion}

Medical world modeling is not a single architecture or application. It is a set of claims about representing state, predicting change, responding to action, comparing alternatives, and using simulated futures to choose decisions. A capability-first structure makes these claims comparable, while SATO-V shows where their evidence is strong or incomplete. The resulting picture is neither that the field is limited to static prediction nor that clinical planning has been solved. Future prediction and action-conditioned simulation are now visible across imaging, EHR, disease progression, physiology, and embodied environments; credible counterfactual reasoning and decision-level validation remain much less integrated with these systems.

The central research task is therefore not generation alone. It is to align a clinically meaningful state with an executable action, an action-sensitive transition, a decision-relevant outcome, and validation appropriate to the intended horizon. This alignment gives lower-level work a precise and useful role, prevents planning language from substituting for evidence, and defines concrete thresholds for progress toward clinically credible world models.
